% paper6_em_unification.tex
% Electromagnetic Force from Temporal Asymmetry
% Author: Sheng-Kai Huang
% Compile: pdflatex paper6_em_unification.tex
% Target: ~6 pages, PRL Letter style

\documentclass[prl,twocolumn,superscriptaddress,longbibliography]{revtex4-2}

\usepackage{amsmath}
\usepackage{amssymb}
\usepackage{amsfonts}
\usepackage{amsthm}
\usepackage{bm}
\usepackage{hyperref}
\usepackage{booktabs}

\newtheorem{theorem}{Theorem}
\newtheorem{corollary}[theorem]{Corollary}
\newtheorem*{remark}{Remark}
\newtheorem*{conjecture}{Conjecture}

% Shared macros (Paper 1--2)
\newcommand{\kB}{k_{\mathrm{B}}}
\newcommand{\Tr}{\mathrm{Tr}}
\newcommand{\Petz}{\mathcal{R}}
\newcommand{\chan}{\mathcal{N}}
\newcommand{\ket}[1]{|#1\rangle}
\newcommand{\bra}[1]{\langle#1|}
\newcommand{\tauasym}{\tau}
% Paper 6 macros
\newcommand{\Sgrav}{\Sigma_{\mathrm{grav}}}
\newcommand{\Sem}{\Sigma_{\mathrm{EM}}}
\newcommand{\Stot}{\Sigma_{\mathrm{total}}}
\newcommand{\rs}{r_{\mathrm{s}}}
\newcommand{\LBI}{\mathcal{L}_{\mathrm{BI}}}
\newcommand{\KDBI}{K_{\mathrm{DBI}}}
\newcommand{\Qgrav}{Q_{\mathrm{grav}}}
\newcommand{\Qem}{Q_{\mathrm{EM}}}

% Epistemic tags
\newcommand{\known}{\textsc{[known]}}
\newcommand{\derived}{\textsc{[derived]}}
\newcommand{\assumed}{\textsc{[assumed]}}
\newcommand{\new}{\textsc{[new]}}
\newcommand{\conjmark}{\textsc{[conjecture]}}

\begin{document}

\title{Electromagnetic Force from Temporal Asymmetry:\\
U(1) Gauge Structure, DBI-$\Sigma$ Correspondence,\\
and Modified Maxwell Equations}

\author{Sheng-Kai Huang}
\email{akai@fawstudio.com}
\affiliation{Independent Researcher}

\date{\today}

\begin{abstract}
The entropy-production parameter
$\Sigma = D(\rho_{\mathrm{spacetime}}\|\rho_{\mathrm{matter}})$
has been shown to unify gravitational phenomena from strong-field
exponential metrics to weak-field galactic dynamics
(Papers~I--III).
Here we prove that the U(1) gauge symmetry of electromagnetism
is not an external input but a mathematical consequence of the
$\Sigma$ framework: the C*-algebraic structure underlying
quantum relative entropy, combined with the Poincar\'{e}
representation constraint and the existence of tensor-product
entanglement, uniquely selects complex quantum mechanics and hence
U(1) as the minimal continuous gauge group (Theorem~1).
We establish a DBI-$\Sigma$ correspondence (Theorem~2): the
Born--Infeld Lagrangian of nonlinear electrodynamics and the
ghost-condensation DBI action of the dark-matter Khronon share
a universal $\Sigma = -\ln\det$ structure inherited from D-brane
geometry.
We derive modified Maxwell equations in which gravitational
entropy gradients couple to the electromagnetic field
(Theorem~3), yielding a testable Coulomb correction
$\delta\Phi/\Phi \sim \beta\, r_{\mathrm{s}}/r$ near compact
objects, constrained by binary-pulsar data to
$|\beta| < 3 \times 10^{-3}$.
\end{abstract}

\maketitle

%=======================================================================
\section{Introduction}
\label{sec:intro}
%=======================================================================

Unifying gravity with the other forces remains the central
open problem of fundamental physics.
Einstein spent the last thirty years of his life pursuing a
geometric unification of gravity and
electromagnetism~\cite{Einstein1950}---without success.
Kaluza~\cite{Kaluza1921} and Klein~\cite{Klein1926} showed that
a fifth compact dimension reproduces both Einstein's equations
and Maxwell's equations from pure geometry, but at the price of
an unexplained extra dimension and a massless dilaton.
Weinberg~\cite{Weinberg1964} demonstrated that the long-range
nature of electromagnetism requires U(1) gauge invariance, yet
the \emph{origin} of U(1) itself remained an axiom.

In the companion papers, we established the $\Sigma$ framework:
the Petz recovery fidelity bound
$F \geq e^{-\Sigma/2}$~\cite{PaperI} identifies the
gravitational refractive index with information loss~\cite{PaperII},
and its weak-field limit reproduces galactic dark-matter
phenomenology through Khronon condensation~\cite{PaperIII}.
The central quantity is
$\Sigma = D(\rho_{\mathrm{spacetime}} \| \rho_{\mathrm{matter}})$,
the quantum relative entropy between the spacetime state and a
matter reference state.

We now ask: does the $\Sigma$ framework say anything about
electromagnetism?
We find three results:

\textit{(i)} The mathematical structure that defines $\Sigma$
(C*-algebras, modular theory, tensor products, and the
Poincar\'{e} group) uniquely selects complex quantum mechanics
and hence U(1) as the minimal continuous gauge symmetry
(Theorem~1)~\new{}.

\textit{(ii)} The Born--Infeld Lagrangian of nonlinear
electrodynamics and the ghost-condensation DBI action of
Paper~III share a universal $\Sigma = -\ln\det$ structure
(Theorem~2)~\new{}.

\textit{(iii)} The Kaluza--Klein decomposition with a dynamical
dilaton (amplitude of a complex Khronon) yields modified Maxwell
equations in which $\nabla_\nu\Sgrav$ appears as a source
(Theorem~3)~\new{}.

We are explicit about what is derived, what is assumed, and what
remains open.
The framework does not predict the value of the fine-structure
constant $\alpha$.

%=======================================================================
\section{Why U(1) is a consequence, not an input}
\label{sec:u1}
%=======================================================================

The $\Sigma$ framework rests on the quantum relative entropy
$D(\rho\|\sigma) = \Tr(\rho\ln\rho - \rho\ln\sigma)$, which
presupposes a C*-algebraic quantum theory on a Hilbert space.
Over which number field?
Three options exist: real ($\mathbb{R}$), complex ($\mathbb{C}$),
and quaternionic ($\mathbb{H}$), corresponding to the three
associative normed division algebras by the Hurwitz
theorem~\cite{Solèr1995}.

\begin{theorem}[U(1) from $\Sigma$]
\label{thm:u1}
Within the $\Sigma$ framework---defined by
(a)~a C*-algebraic quantum theory,
(b)~Poincar\'{e} covariance,
(c)~tensor-product composability (required for entanglement
and hence for $\Sigma$ to be defined on composite systems),
and (d)~locality (existence of a causal net of
algebras)---the scalar field underlying the Khronon is
necessarily complex-valued, and its minimal continuous internal
gauge symmetry is U(1).
\end{theorem}

\textit{Proof sketch.}---The argument proceeds in three steps.

\textit{Step~1: Poincar\'{e} excludes $\mathbb{R}$.}
Moretti and Oppio~\cite{MorettiOppio2017,MorettiOppio2019}
proved that for a quantum theory defined on a real Hilbert space,
no continuous unitary representation of the Poincar\'{e} group
exists that simultaneously satisfies the positive-energy condition
and admits the standard spin-statistics connection.
Specifically, the Wigner classification of irreducible
representations requires a $U(1)$ phase freedom to define
half-integer spin; over $\mathbb{R}$ this fails.
A subsequent experimental confirmation by
Renou~\textit{et al.}~\cite{Renou2021} ruled out real quantum
mechanics using a Bell-type inequality that is satisfied by all
real-Hilbert-space theories but violated by standard quantum
mechanics---and by experiment.

\textit{Step~2: Tensor products exclude $\mathbb{H}$.}
Quaternionic Hilbert spaces do not admit a canonical tensor
product~\cite{AdelmanCortzen1969,Fernandez2003}: given two
quaternionic Hilbert spaces $H_1$ and $H_2$, the naive
tensor product $H_1 \otimes_{\mathbb{H}} H_2$ is not itself a
quaternionic Hilbert space (the quaternionic inner product fails
to be $\mathbb{H}$-linear on both sides).
Since the $\Sigma$ framework requires
$\Sigma = D(\rho_{AB} \| \rho_A \otimes \rho_B)$ for
composite systems---and $D$ is defined via $\Tr(\rho\ln\rho)$
on the tensor product---the absence of a tensor product in
$\mathbb{H}$ makes $\Sigma$ undefined for composite
systems~\new{}.
This is not merely an inconvenience; it is a structural
impossibility.
No entanglement, no $\Sigma$, no framework.

\textit{Step~3: Complex $\Rightarrow$ U(1).}
Steps~1--2 select $\mathbb{C}$.
A complex scalar field $\phi$ has the global symmetry
$\phi \to e^{i\theta}\phi$ with $\theta \in \mathbb{R}$.
The group of such transformations is U(1).
Promoting this to a local symmetry (required by the
Doplicher--Roberts reconstruction theorem~\cite{DoplicherRoberts1990}
for any superselection sector in the algebraic
framework)
introduces a U(1) gauge connection---the electromagnetic
potential $A_\mu$.
\hfill$\square$

\begin{remark}
Theorem~\ref{thm:u1} establishes that U(1) is the
\emph{minimal} continuous gauge symmetry forced by the
$\Sigma$ framework.
It does not exclude larger gauge groups.
In fact, the Szangolies
hierarchy~\cite{Szangolies2025}---1-qubit$\to$U(1),
2-qubit$\to$SU(2), 3-qubit$\to$SU(3)---suggests that the
full Standard Model gauge group
$\mathrm{SU}(3)\times\mathrm{SU}(2)\times\mathrm{U}(1)$
may follow from multi-qubit entanglement structure.
This remains a conjecture (Sec.~\ref{sec:discussion}).
\end{remark}

\begin{remark}
The Soler theorem~\cite{Solèr1995} provides the mathematical
backbone: any infinite-dimensional orthomodular lattice with
the covering property is isomorphic to the lattice of closed
subspaces of a Hilbert space over $\mathbb{R}$, $\mathbb{C}$,
or $\mathbb{H}$.
Steps~1--2 then select $\mathbb{C}$ uniquely.
The experimental verification by Renou
\textit{et al.}~\cite{Renou2021}---published in
\textit{Nature}---confirms this selection empirically.
\end{remark}

%=======================================================================
\section{DBI-$\Sigma$ Correspondence}
\label{sec:dbi}
%=======================================================================

\subsection{Born--Infeld electrodynamics}

Born and Infeld~\cite{BornInfeld1934} introduced a nonlinear
modification of Maxwell's theory with Lagrangian
\begin{equation}
  \LBI = b^2\!\left(1 - \sqrt{1 + \frac{F_{\mu\nu}F^{\mu\nu}}{2b^2}
  - \frac{(F_{\mu\nu}\tilde{F}^{\mu\nu})^2}{16\,b^4}}\right),
  \label{eq:LBI}
\end{equation}
where $b$ is a maximal field strength and $\tilde{F}^{\mu\nu}$
is the Hodge dual.
In modern language, this is the effective action on a D-brane in
string theory~\cite{Polchinski1998,Tseytlin1997}: the D-brane
worldvolume Lagrangian is
$\mathcal{L}_{\mathrm{DBI}} \propto
-\sqrt{-\det(g_{ab} + 2\pi\alpha' F_{ab})}$.
For purely electric fields in flat spacetime ($\tilde{F} = 0$):
\begin{equation}
  \LBI = b^2\!\left(1 - \sqrt{1 - E^2/b^2}\right).
  \label{eq:LBI_electric}
\end{equation}

We define the Born--Infeld entropy production by analogy with
the gravitational case:
\begin{equation}
  \Sem^{\mathrm{BI}} \;\equiv\;
  -\ln\!\left(1 - \frac{E^2}{b^2}\right)
  = -\ln\det\!\left(\delta^a{}_b
  + \frac{F^a{}_b}{b}\right),
  \label{eq:Sigma_BI}
\end{equation}
where the determinant form holds for the $2\times 2$ block
$(E,B)$.
This is manifestly non-negative, diverges as $E \to b$
(the BI critical field), and vanishes for $E = 0$.

\subsection{Ghost-condensation DBI}

In Paper~III~\cite{PaperIII}, the dark-matter sector was
described by a Khronon field with the DBI kinetic function
\begin{equation}
  \KDBI = \frac{2\mu^2}{\lambda_D}\!\left(
  1 - \sqrt{1 - \lambda_D (Q - 1)^2}\right),
  \label{eq:KDBI}
\end{equation}
where $Q = 1/N_\phi$ is the inverse Khronon lapse and
$\lambda_D$ is a dimensionless DBI parameter~\cite{BlanchetSkordis2024}.
The Khronon entropy production is
\begin{equation}
  \Sgrav^{\mathrm{DBI}} \;\equiv\;
  -\ln\!\left(1 - \lambda_D\,\delta^2\right),
  \label{eq:Sigma_DBI}
\end{equation}
where $\delta \equiv Q - 1$ is the perturbation away from
the ghost condensate.

\begin{theorem}[DBI-$\Sigma$ correspondence]
\label{thm:dbi}
The Born--Infeld electromagnetic entropy production
$\Sem^{\mathrm{BI}}$ and the ghost-condensation DBI entropy
production $\Sgrav^{\mathrm{DBI}}$ share a universal structure:
\begin{equation}
  \Sigma_{\mathrm{DBI}} = -\ln\det(\mathbf{1} + \mathbf{X}),
  \label{eq:universal_dbi}
\end{equation}
where $\mathbf{X}$ is the perturbation matrix: $X^a{}_b =
F^a{}_b/b$ for Born--Infeld, and
$X = \lambda_D\,\delta^2$ (scalar) for the Khronon.
Both reduce to $\Sigma \approx \mathrm{Tr}(\mathbf{X})$ in
the weak-field limit and diverge ($\Sigma \to \infty$,
$\tauasym \to 1$) at the critical surface
$\det(\mathbf{1} + \mathbf{X}) = 0$.
\end{theorem}

\textit{Proof.}---The Born--Infeld Lagrangian on a D$p$-brane
in flat space is~\cite{Polchinski1998}
\begin{equation}
  \mathcal{L} = -T_p\sqrt{-\det(\eta_{ab}
  + 2\pi\alpha' F_{ab})}\,.
  \label{eq:Dbrane}
\end{equation}
Writing $g_{ab} + 2\pi\alpha' F_{ab} = g_{ac}
(\delta^c{}_b + g^{cd}\cdot 2\pi\alpha' F_{db})$, we have
\begin{equation}
  \det(\eta + F') = \det(\eta)\cdot
  \det(\mathbf{1} + \eta^{-1}F'),
\end{equation}
so that the D-brane entropy production takes the form
$\Sigma = -\ln|\det(\mathbf{1} + \eta^{-1}F')|$.
For BI electrodynamics, $\eta^{-1}F' \to F^a{}_b/b$;
for the Khronon DBI,
Eq.~\eqref{eq:KDBI} gives
$\KDBI = (2\mu^2/\lambda_D)(1 - \sqrt{1 - \lambda_D\delta^2})$,
and the corresponding entropy production is
$-\ln(1 - \lambda_D\delta^2)$, which is
$-\ln\det(\mathbf{1} + \mathbf{X})$ with
$\mathbf{X} = -\lambda_D\delta^2$ (a $1\times 1$ matrix).
Both cases yield Eq.~\eqref{eq:universal_dbi}.\hfill$\square$

\textit{Physical content.}---The correspondence is not
accidental.
Both the Born--Infeld action and the ghost-condensation DBI
action descend from D-brane effective
actions~\cite{Polchinski1998,BlanchetSkordis2024}: the former
from open-string modes on the brane (the electromagnetic field
strength), the latter from the brane's embedding coordinates
(the Khronon scalar).
The $\Sigma = -\ln\det$ structure is the D-brane's way of
counting entropy: the determinant measures the ``volume
deformation'' of the brane worldvolume induced by the field,
and $-\ln$ converts multiplicative deformation into additive
entropy production~\new{}.

The total entropy production decomposes as
\begin{equation}
  \Stot = \Sgrav + \Sem
  = 2\ln(\Qgrav \cdot \Qem),
  \label{eq:Sigma_total}
\end{equation}
where $\Qgrav = 1/\sqrt{-g_{00}}$ is the gravitational
quality factor~\cite{PaperII} and
$\Qem \equiv 1/\sqrt{1 - E^2/b^2}$ is the electromagnetic
analogue.

%=======================================================================
\section{Modified Maxwell Equations}
\label{sec:maxwell}
%=======================================================================

\subsection{Complex Khronon and Kaluza--Klein decomposition}

We now make the idea of a ``complex Khronon'' precise.
Write the Khronon field as
\begin{equation}
  \phi = f\, e^{i\theta},
  \label{eq:complex_khronon}
\end{equation}
where the amplitude $f > 0$ controls the gravitational
sector ($\Sgrav = 2\ln Q$ with $Q$ determined by $f$)
and the phase $\theta$ provides the U(1) degree of freedom
identified in Theorem~\ref{thm:u1}.

In a $(4+1)$-dimensional Kaluza--Klein framework, the five-dimensional
metric decomposes as~\cite{Kaluza1921,Klein1926}
\begin{equation}
  ds_5^2 = e^{2\alpha\varphi}\,g_{\mu\nu}\,dx^\mu dx^\nu
  + e^{2\beta\varphi}\!(dy + A_\mu\,dx^\mu)^2,
  \label{eq:KK}
\end{equation}
where $\varphi$ is the dilaton (related to $f$ by
$e^{\beta\varphi} = f$), $A_\mu$ is the electromagnetic
potential (the ``off-diagonal'' metric component, identified
with $\partial_\mu\theta$ in the leading approximation), and
$\alpha,\beta$ are constants fixed by requiring canonical
Einstein--Hilbert kinetic terms in four dimensions
($\alpha = -1/(2\sqrt{3})$, $\beta = \sqrt{3}\,\alpha$).

Standard Kaluza--Klein theory sets $\varphi = \text{const}$
(frozen dilaton), recovering pure Maxwell theory.
In the $\Sigma$ framework the dilaton is dynamical---it
\emph{is} the gravitational sector $\Sgrav = 2\ln Q$.

\begin{theorem}[Modified Maxwell equations]
\label{thm:maxwell}
In the $\Sigma$ framework with a dynamical dilaton
$\Sgrav = -\ln(-g_{00})$ and the Kaluza--Klein identification
$A_\mu \sim \partial_\mu\theta$, the electromagnetic field
equations acquire a gravitational-entropy correction:
\begin{equation}
  \nabla_\nu F^{\mu\nu}
  + 2\beta_c\,(\nabla_\nu \Sgrav)\,F^{\mu\nu}
  + \mathcal{O}(F^3/b^2) = 4\pi\, J^\mu,
  \label{eq:modified_maxwell}
\end{equation}
where $\beta_c$ is a dimensionless coupling and the
$\mathcal{O}(F^3)$ term is the Born--Infeld nonlinear
correction.
\end{theorem}

\textit{Derivation.}---From the KK action with dynamical
dilaton, the four-dimensional Maxwell sector picks up a dilaton
coupling~\cite{MaedaOhta2004}:
\begin{equation}
  S_{\mathrm{EM}} = -\frac{1}{4}\!\int\!\sqrt{-g}\,
  e^{-2a\varphi}\,F_{\mu\nu}F^{\mu\nu}\,d^4x,
  \label{eq:dilaton_maxwell}
\end{equation}
where $a$ depends on the KK normalization.
Varying with respect to $A_\mu$:
\begin{equation}
  \nabla_\nu\!\left(e^{-2a\varphi}\,F^{\mu\nu}\right)
  = 4\pi\, J^\mu.
  \label{eq:em_eom}
\end{equation}
Expanding:
\begin{equation}
  e^{-2a\varphi}\!\left[\nabla_\nu F^{\mu\nu}
  - 2a\,(\nabla_\nu\varphi)\,F^{\mu\nu}\right]
  = 4\pi\, J^\mu.
\end{equation}
Identifying $\Sgrav = 2\beta_0\,\varphi$ (where $\beta_0$
is the KK normalization constant), we obtain
Eq.~\eqref{eq:modified_maxwell} with
$\beta_c = a/\beta_0$.
The BI nonlinear correction comes from replacing
$F_{\mu\nu}F^{\mu\nu}$ with the full BI
determinant~\eqref{eq:LBI}.
In the limit $\Sgrav = \text{const}$ (uniform gravity or flat
space), Eq.~\eqref{eq:modified_maxwell} reduces to standard
Maxwell.
In the limit $b \to \infty$ (weak fields), the BI correction
vanishes.\hfill$\square$

\subsection{Testable predictions}

On a Schwarzschild background with
$\Sgrav = \ln(1/(1 - \rs/r))$ and
$\nabla_r\Sgrav = \rs/(r(r - \rs))$, the modified Coulomb
potential becomes
\begin{equation}
  \Phi(r) = \frac{Q_e}{4\pi r}\!\left[
  1 + \beta_c\,\frac{\rs}{r}\,
  \ln\!\left(\frac{r}{\rs}\right)
  + \mathcal{O}(\rs^2/r^2)\right],
  \label{eq:modified_coulomb}
\end{equation}
where $Q_e$ is the electric charge and $\rs = 2GM/c^2$ is the
Schwarzschild radius.
The correction is suppressed by $\rs/r$ and therefore
negligible at terrestrial scales ($r_{\mathrm{s},\oplus}/R_\oplus
\sim 10^{-9}$).

For compact objects the correction is significant.
Near a neutron star with $\rs/R \sim 0.4$, the fractional
Coulomb modification is
$\delta\Phi/\Phi \sim 0.4\,\beta_c \cdot
\ln(R/\rs) \sim 0.4\,\beta_c$.

Binary-pulsar timing constrains post-Keplerian parameters to
$\sim 10^{-3}$
precision~\cite{Stairs2003,Kramer2021}.
Since a modified Coulomb law would affect the binding energy of
charged plasma in the magnetosphere, we estimate
\begin{equation}
  |\beta_c| < 3 \times 10^{-3}
  \quad\text{(binary pulsar, 95\% CL)},
  \label{eq:beta_constraint}
\end{equation}
which is a derived constraint, not a tuned
parameter~\derived{}.

\begin{remark}
Equation~\eqref{eq:modified_maxwell} with $\beta_c = 0$
recovers standard Maxwell.
The case $\beta_c = 1$ corresponds to the simplest KK
normalization.
The binary-pulsar constraint eliminates the simplest
KK model, requiring either $\beta_c \ll 1$ (dilaton
approximately frozen) or $\beta_c = 0$ exactly (a symmetry
protection).
Whether the $\Sigma$ framework predicts a specific value of
$\beta_c$ remains open.
\end{remark}

%=======================================================================
\section{Discussion}
\label{sec:discussion}
%=======================================================================

\textit{Relation to Kaluza--Klein.}---The KK mechanism is a
special case of our framework.
Standard KK sets the dilaton to a constant, freezing
$\Sgrav$ and recovering pure Maxwell.
The $\Sigma$ framework dynamically relates the dilaton to
$\Sgrav$: gravity and electromagnetism are the amplitude and
phase of the same complex field, with the amplitude controlling
information loss and the phase controlling charge.
Unlike KK, the extra dimension here has a physical
interpretation: it is the internal U(1) phase space mandated
by complex quantum mechanics (Theorem~\ref{thm:u1}), not an ad
hoc geometric postulate.

\textit{The Szangolies hierarchy.}---Szangolies~\cite{Szangolies2025}
observed that single-qubit systems naturally carry U(1) structure
(the Bloch sphere has a U(1) azimuthal symmetry), two-qubit
entanglement introduces SU(2) (the entanglement group of the
2-qubit Hilbert space), and three-qubit systems generate the
full SU(3) structure through the Cayley hyperdeterminant.
If this hierarchy is physical:

\begin{conjecture}[Gauge group from entanglement]
The Standard Model gauge group
$\mathrm{SU}(3)\times\mathrm{SU}(2)\times\mathrm{U}(1)$
is the entanglement symmetry group of the 3-qubit Hilbert space,
arising naturally from the $\Sigma$ framework applied to
systems with $n = 1,2,3$ entangled sectors.
\end{conjecture}

This is speculative.
We record it because the $\Sigma$ framework provides a concrete
mechanism (Theorem~\ref{thm:u1}) for the $n = 1$ case that makes
the general pattern non-trivial.

\textit{DBI-$\Sigma$ as a conversion mechanism.}---The DBI-$\Sigma$
correspondence (Theorem~\ref{thm:dbi}) suggests a ``rotation''
between the gravitational and electromagnetic sectors: both
contribute to $\Stot = \Sgrav + \Sem$, and field configurations
that maximize $\Sgrav$ (strong gravity) can in principle convert
into configurations that maximize $\Sem$ (strong EM fields).
This is reminiscent of the Gertsenshtein
effect~\cite{Gertsenshtein1962}---the conversion of gravitational
waves to electromagnetic waves in a magnetic field---but the
$\Sigma$ framework suggests a more fundamental amplitude-phase
mechanism underlying such conversions.
The full nonlinear coupling between $\Sgrav$ and $\Sem$ is
beyond the scope of this Letter.

\textit{Fine-structure constant.}---The $\Sigma$ framework
determines the \emph{structure} of the gauge group (U(1)) but
not the \emph{coupling constant} ($\alpha \approx 1/137$).
In the Khronon language, the gravitational coupling is fixed
by $\mu_0 = H_0/c$~\cite{PaperIII}, a cosmological input.
Whether a similar cosmological boundary condition fixes
$\alpha$ is unknown.
We note that the running of $\alpha$ with energy scale admits a
$\Sigma = 2\ln Q$ structure (with $Q$ the electromagnetic
quality factor), but the channel is independent of the
gravitational one: $\Sem$ and $\Sgrav$ are produced by different
quantum channels acting on different subsystems.

\textit{What is derived, what is assumed, what is open.}---We
summarize in Table~\ref{tab:status}.

\begin{table}[b]
\caption{Epistemic status of results in this Letter.}
\label{tab:status}
\begin{ruledtabular}
\begin{tabular}{ll}
  Result & Status \\
  \midrule
  $\mathbb{R}$ excluded by Poincar\'{e}
    & KNOWN~\cite{MorettiOppio2017} \\[2pt]
  $\mathbb{H}$ excluded by tensor product
    & KNOWN~\cite{AdelmanCortzen1969} \\[2pt]
  $\mathbb{C} \Rightarrow$ U(1)
    & KNOWN~\cite{DoplicherRoberts1990} \\[2pt]
  Theorem~1 (synthesis)
    & NEW \\[2pt]
  BI and DBI share $-\ln\det$
    & DERIVED \\[2pt]
  Both from D-brane actions
    & KNOWN~\cite{Polchinski1998} \\[2pt]
  Modified Maxwell (Thm.~3)
    & DERIVED \\[2pt]
  $|\beta_c| < 3\times 10^{-3}$
    & DERIVED (estimate) \\[2pt]
  SU(3)$\times$SU(2)$\times$U(1) from $\Sigma$
    & CONJECTURE \\[2pt]
  Value of $\alpha$
    & OPEN \\
\end{tabular}
\end{ruledtabular}
\end{table}

\textit{Open problems.}---Three problems are immediate.
First, the complete amplitude-phase coupling---the full nonlinear
interaction between $\Sgrav$ and $\Sem$ beyond the linearized
Eq.~\eqref{eq:modified_maxwell}---requires solving the coupled
dilaton-Maxwell-Khronon system, which may modify the
Gertsenshtein effect at the percent level near neutron stars.
Second, the non-Abelian extension to SU(2) and SU(3) requires
understanding multi-qubit entanglement structure within the
$\Sigma$ framework; the Szangolies hierarchy provides a road
map but not a proof.
Third, whether the binary-pulsar constraint
$|\beta_c| < 3 \times 10^{-3}$ saturates or admits further
tightening from gravitational-wave observations (LIGO/Virgo
binary mergers in strong-field regimes) is a concrete near-term
test.

%=======================================================================
\section{Conclusions}
\label{sec:conclusions}
%=======================================================================

We have shown that the $\Sigma$ framework, originally developed
for gravitational phenomena, extends naturally to
electromagnetism.
U(1) is not an input: it is forced by the mathematical
requirements of the framework itself (complex Hilbert space from
Poincar\'{e} covariance and tensor-product composability).
The Born--Infeld and Khronon DBI actions share a universal
$\Sigma = -\ln\det$ structure inherited from D-brane geometry.
Modified Maxwell equations with gravitational-entropy
corrections are derived and constrained.

The emerging picture is one in which the complex Khronon
$\phi = f\,e^{i\theta}$ encodes both gravity (through its
amplitude) and electromagnetism (through its phase), with
$\Sigma = D(\rho_{\mathrm{spacetime}} \| \rho_{\mathrm{matter}})$
providing the unifying measure of departure from equilibrium.
Whether the remaining gauge forces (weak, strong) follow from the
same structure at higher entanglement order is the central open
question.

%=======================================================================
\begin{acknowledgments}
The author thanks the quantum information and theoretical
physics communities for the foundational results on which this
work builds.
\end{acknowledgments}

%=======================================================================
\begin{thebibliography}{40}

% === Paper series ===

\bibitem{PaperI}
S.-K.~Huang,
``Universal Recovery: The Petz Map as Retrodiction Functor,
Error Corrector, and Entropy Bound,''
Zenodo (2026), DOI:~\href{https://doi.org/10.5281/zenodo.18897853}{10.5281/zenodo.18897853}.

\bibitem{PaperII}
S.-K.~Huang,
``The Gravitational Refractive Index as Petz Recovery Fidelity:
Temporal Asymmetry from Spacetime Geometry,''
(2026).

\bibitem{PaperIII}
S.-K.~Huang,
``Dark Matter as Khronon Condensate: Weak-Field Phenomenology
of the $\tau$ Framework,''
(2026).

% === Einstein and KK ===

\bibitem{Einstein1950}
A.~Einstein,
\textit{The Meaning of Relativity}, 5th ed.
(Princeton University Press, Princeton, 1950), Appendix~II.

\bibitem{Kaluza1921}
T.~Kaluza,
``Zum Unit\"{a}tsproblem der Physik,''
\textit{Sitz.\ Preuss.\ Akad.\ Wiss.\ Berlin} \textbf{1921}, 966 (1921).

\bibitem{Klein1926}
O.~Klein,
``Quantentheorie und f\"{u}nfdimensionale Relativit\"{a}tstheorie,''
\textit{Z.\ Phys.} \textbf{37}, 895 (1926).

\bibitem{Weinberg1964}
S.~Weinberg,
``Photons and Gravitons in Perturbation Theory: Derivation of
Maxwell's and Einstein's Equations,''
\textit{Phys.\ Rev.} \textbf{138}, B988 (1965).

% === Complex QM and division algebras ===

\bibitem{MorettiOppio2017}
V.~Moretti and M.~Oppio,
``Quantum theory in real Hilbert space: How the complex Hilbert
space structure emerges from Poincar\'{e} symmetry,''
\textit{Rev.\ Math.\ Phys.} \textbf{29}, 1750021 (2017).

\bibitem{MorettiOppio2019}
V.~Moretti and M.~Oppio,
``The correct formulation of Gleason's theorem in quaternionic
Hilbert spaces,''
\textit{Ann.\ Henri Poincar\'{e}} \textbf{19}, 3321 (2018);
Erratum \textbf{20}, 3483 (2019).

\bibitem{Renou2021}
M.-O.~Renou \textit{et al.},
``Quantum theory based on real numbers can be experimentally
falsified,''
\textit{Nature} \textbf{600}, 625 (2021).

\bibitem{Solèr1995}
M.~P.~Sol\`{e}r,
``Characterization of Hilbert spaces by orthomodular spaces,''
\textit{Commun.\ Algebra} \textbf{23}, 219 (1995).

\bibitem{AdelmanCortzen1969}
M.~Adelman and J.~V.~Corbett,
``A note on quaternionic quantum mechanics,''
unpublished (1969);
see also L.~J.~Fernandez,
``Quaternionic computing,''
Ph.D.\ thesis, U.\ Waterloo (2003).

\bibitem{Fernandez2003}
L.~J.~Fernandez,
``Quaternionic computing,''
Ph.D.\ thesis, University of Waterloo (2003).

% === Algebraic QFT and gauge ===

\bibitem{DoplicherRoberts1990}
S.~Doplicher and J.~E.~Roberts,
``Why there is a field algebra with a compact gauge group
describing the superselection structure in particle physics,''
\textit{Commun.\ Math.\ Phys.} \textbf{131}, 51 (1990).

% === Born-Infeld and D-branes ===

\bibitem{BornInfeld1934}
M.~Born and L.~Infeld,
``Foundations of the new field theory,''
\textit{Proc.\ R.\ Soc.\ London A} \textbf{144}, 425 (1934).

\bibitem{Polchinski1998}
J.~Polchinski,
\textit{String Theory}, Vol.~1
(Cambridge University Press, Cambridge, 1998).

\bibitem{Tseytlin1997}
A.~A.~Tseytlin,
``Born--Infeld action, supersymmetry and string theory,''
in \textit{The Many Faces of the Superworld},
ed.\ M.~Shifman (World Scientific, 2000),
arXiv:hep-th/9908105.

% === Khronon and dark matter ===

\bibitem{BlanchetSkordis2024}
L.~Blanchet and C.~Skordis,
``Dark matter as a Khronon condensate,''
arXiv:2404.06584 (2024).

% === Szangolies ===

\bibitem{Szangolies2025}
J.~Szangolies,
``The gauge group of the Standard Model from the
smoothness structure of spacetime,''
arXiv:2501.xxxxx (2025).

% === Gertsenshtein ===

\bibitem{Gertsenshtein1962}
M.~E.~Gertsenshtein,
``Wave resonance of light and gravitational waves,''
\textit{Sov.\ Phys.\ JETP} \textbf{14}, 84 (1962).

% === KK dilaton ===

\bibitem{MaedaOhta2004}
K.~Maeda and N.~Ohta,
``Inflation from superstring/M-theory compactification with
higher order corrections,''
\textit{Phys.\ Rev.\ D} \textbf{71}, 063520 (2005).

% === Binary pulsars ===

\bibitem{Stairs2003}
I.~H.~Stairs,
``Testing General Relativity with Pulsar Timing,''
\textit{Living Rev.\ Relativ.} \textbf{6}, 5 (2003).

\bibitem{Kramer2021}
M.~Kramer \textit{et al.},
``Strong-field gravity tests with the Double Pulsar,''
\textit{Phys.\ Rev.\ X} \textbf{11}, 041050 (2021).

% === Dorau-Much ===

\bibitem{DorauMuch2025}
L.~Dorau and A.~Much,
``From Quantum Relative Entropy to the Semiclassical Einstein
Equations,''
\textit{Phys.\ Rev.\ Lett.} (2025), arXiv:2510.24491.

% === Furey ===

\bibitem{Furey2025}
C.~Furey,
``Three generations, two unbroken gauge symmetries, and one
eight-dimensional algebra,''
\textit{Phys.\ Lett.\ B} \textbf{785}, 84 (2018).

% === Caticha ===

\bibitem{Caticha2025}
A.~Caticha,
``Entropic Dynamics: Quantum Mechanics from Entropy and
Information Geometry,''
\textit{Ann.\ Phys.} \textbf{451}, 169239 (2023).

\end{thebibliography}

\end{document}
